\section{Rectangular Waveguide Theory}
\label{sec:rectangular-waveguide-theory}

For a vacuum-filled rectangular waveguide, the cutoff frequency of a TE or TM mode is \cite{pozar2011microwave}

\[
f_{c,mn}=
\frac{c}{2}
\sqrt{
\left(\frac{m}{a}\right)^2+
\left(\frac{n}{b}\right)^2
}.
\]

The transverse cutoff wavenumber and longitudinal propagation constant are

\[
k_c=
\sqrt{
\left(\frac{m\pi}{a}\right)^2+
\left(\frac{n\pi}{b}\right)^2
},
\qquad
\beta_{mn}=\sqrt{k_0^2-k_c^2},
\]

where \(k_0=2\pi f/c\). A mode is propagating when \(f>f_c\) and evanescent when \(f<f_c\).

\subsection{Lowest-Order Cutoff Frequencies}

For \(a=b=20~\mathrm{mm}\),

\[
f_{c,10}=f_{c,01}=\frac{c}{2a}=7.49481145~\mathrm{GHz}.
\]

The next cutoff family is

\[
f_{c,11}
=
\frac{c}{2}
\sqrt{\frac{1}{a^2}+\frac{1}{b^2}}
=
10.599264~\mathrm{GHz}.
\]

The \(\mathrm{TE}_{20}\) and \(\mathrm{TE}_{02}\) modes begin near

\[
f_{c,20}=f_{c,02}=\frac{c}{a}=14.989623~\mathrm{GHz}.
\]

\begin{table}[H]
\centering
\caption{Analytical low-order modal behavior at \(10~\mathrm{GHz}\).}
\label{tab:analytical-modes}
\begin{tabular}{@{}llll@{}}
\toprule
\textbf{Mode family} & \textbf{Cutoff} & \textbf{Status at 10 GHz} & \textbf{Comment} \\
\midrule
\(\mathrm{TE}_{10}\) & \(7.4948~\mathrm{GHz}\) & Propagating & Degenerate lowest-order pair \\
\(\mathrm{TE}_{01}\) & \(7.4948~\mathrm{GHz}\) & Propagating & Orthogonal polarization partner \\
\((1,1)\) TE/TM family & \(10.5993~\mathrm{GHz}\) & Evanescent & Slightly below cutoff \\
\(\mathrm{TE}_{20}/\mathrm{TE}_{02}\) & \(14.9896~\mathrm{GHz}\) & Evanescent & Near upper sweep boundary \\
\bottomrule
\end{tabular}
\end{table}

\subsection{Guide Wavelength and Electrical Length}

For a propagating mode,

\[
\lambda_g
=
\frac{\lambda_0}{
\sqrt{1-\left(f_c/f\right)^2}
},
\qquad
\lambda_0=\frac{c}{f}.
\]

At \(10~\mathrm{GHz}\),

\[
\lambda_0=29.979~\mathrm{mm},
\]

and for either member of the lowest-order pair,

\[
\lambda_g=45.284~\mathrm{mm}.
\]

The \(90~\mathrm{mm}\) physical section is therefore

\[
\frac{l}{\lambda_g}=1.987
\]

guide wavelengths long.

\subsection{TE Wave Impedance}

The TE-mode wave impedance is

\[
Z_{\mathrm{TE}}
=
\frac{\eta_0}{
\sqrt{1-\left(f_c/f\right)^2}
}.
\]

Using \(\eta_0=376.730~\Omega\), the lowest-order TE impedance at \(10~\mathrm{GHz}\) is

\[
Z_{\mathrm{TE}}\approx569.06~\Omega.
\]

This modal impedance is not a fixed \(50~\Omega\) normalization. A waveguide port solves the transverse eigenfield and computes its frequency-dependent propagation and impedance characteristics \cite{simulia_learning_tutorials,simulia_rf_components}.

\subsection{Degenerate Modal Subspaces}

When two modes have the same eigenvalue, any orthogonal linear combination of those modes is also a valid eigenbasis \cite{pozar2011microwave}. For the square guide, a solver may return one pair of orthogonal transverse field patterns at the input and a rotated pair at the output. If \(\mathbf{u}_1,\mathbf{u}_2\) span the input degenerate subspace and \(\mathbf{v}_1,\mathbf{v}_2\) span the output subspace, then

\[
\begin{bmatrix}
\mathbf{v}_1\\
\mathbf{v}_2
\end{bmatrix}
=
\mathbf{R}
\begin{bmatrix}
\mathbf{u}_1\\
\mathbf{u}_2
\end{bmatrix},
\]

where \(\mathbf{R}\) is an orthogonal basis rotation. Consequently, a near-zero \(S_{2(1),1(1)}\) and near-unity \(S_{2(2),1(1)}\) can represent the same physical polarization transported through a uniform guide, merely expressed in a different numerical basis.
